\section{Discussion and limitations}
\label{sec:discussion}

The two remedies make different tradeoffs.  Exact interval bisection followed
by \eqref{eq:KF-main} has no smoothing bias but needs a rank decision.
Pseudo-Huber gives a strictly monotone equation and converges to \(K_R\),
whereas the family in \cref{prop:smoothing-family-main} identifies \(K_F\) as
the Huber Fenchel center.  These centers agree in precisely the cases of
\cref{cor:center-agreement-main}; in particular, they agree for rank-one
normals.  The boundary bias is profile dependent:
polynomial tail order \(p\) gives \(\delta^{p/(p+1)}\), pseudo-Huber gives
\(\delta^{2/3}\), and finite saturation gives \(\delta\).  At fixed strict
crossings the first direction coefficient is computable from
\eqref{eq:frozen-field-main}; compact-root instances instead begin at
cubic/quadratic order and obey the multiplier hierarchy of
\cref{rem:compact-root-hierarchy-main}.

The strict laws allow arbitrary shape and corank, but boundary and crossover
laws currently require a transverse square corank-one crossing; multiple
signed small modes and nontransverse exponents remain open.  The admissible
family does not cover every vanishing regularizer, and full-space openness
implies no frequency law without a distribution for \((G,\Theta)\).  No outer
complexity or learning-performance gain follows without a separate descent
model and control of the changing normal.

Finite Newton--Schulz iterations are separate numerical maps, not exact
polar factors or automatically members of our Fenchel family.
Li and Tsuchiya~\cite{li2026finite} analyze a smoothing benefit of finite iterations for
nonsmooth nonconvex optimization.  Their guarantees use classical
Taylor-derived coefficients and a normalization fixed before the run;
they do not directly cover the tuned quintic with data-dependent
Frobenius normalization used in common Muon/SSO implementations.  Neither
an exact compact-root counterexample nor a small finite-iteration scalar
residual settles the practical method's training performance.  Numerical
validation should measure tangency, norm feasibility, and a valid
primal--dual gap on the same returned direction; estimated-normal and
true-normal residuals should be distinguished.

The experiments in this paper check matrix identities, asymptotic constants,
and certificate diagnostics.  Language-model training would test a separate
application claim.  In particular, reproducing all architectures and training
budgets of Xie et al.~\cite{xie2026controlled} is not required to validate the matrix
analysis results, whereas a claim of matching or improving their reported
training outcomes would require a separately matched protocol.

\section{Conclusion}

We characterized the tangency-constrained certificate slice for nuclear-norm
approximation from a matrix line and identified its hard and Fenchel spectral
centers.  An admissible smoothing family links center selection and boundary
exponents, while strict arbitrary-corank paths admit computable analytic
multiplier and direction expansions.  Transverse corank-one boundaries
exhibit the sharp pseudo-Huber two-thirds law and a uniform cubic crossover.
Compact-polar open failure explains why the full certificate slice is
essential.

\section*{Data and code availability}

The mathematical examples use explicitly specified matrices and no external
training data.  The accompanying source archive contains the numerical
verification scripts, high-precision table generation, and dependency
information.  These computations are reproducibility checks in the recorded
environment, not machine-rigorous interval certificates or formal proofs.
Separately supplied GPU application protocols have no reported training
results in this manuscript.

\section*{Declaration on generative AI assistance}

Generative AI tools assisted with language editing, literature discovery,
algebraic exploration, drafting portions of theorem and proof text, and
drafting verification code.  Responsibility for the mathematical arguments, source attribution,
computational claims, and final submission remains with the human authors.
